% 2026-09-01 exact cross-frontier synthesis: depth-3 regulus -> E8 completion chart.
\subsection{Depth-three reguli as unique $E_8$ completion addresses}
\label{sec:20260901-regulus-e8-completion}

The depth-three blocking computation in the Holotrade track and the selected
$E_8$/$D_4$ geometry in the present repository meet objectwise, not merely
numerically.  The forty-five minimum weight-eight tritangent supports are the
forty-five orthogonal $D_4\oplus D_4$ packets selected by the $W(3,3)$
polarity.  Their disjointness graph is the point graph of $GQ(4,2)$: it has
$45$ vertices, degree $12$, and $270$ edges.  Independently, the Holotrade
blocking audit identifies its $270$ all-isotropic depth-three obstruction
reguli with exactly these $270$ support-disjoint packet pairs.

The twenty-seven maximal five-cliques of this graph give a stronger completion
law.  Every disjoint packet pair is contained in exactly one five-clique, and
every five-clique partitions the forty $W(3,3)$ fibres.  Under the certified
six-roots-per-fibre $E_8$ lift, each packet contains two orthogonal $D_4$
subsystems and therefore $48$ roots.  Consequently
\[
  \boxed{270=27\binom{5}{2}}
\]
is an objectwise incidence identity: each depth-three obstruction regulus
selects two $48$-root packets, hence $96$ roots, and determines a unique
completion by the other three packets, hence another $144$ roots, to one of
twenty-seven partitions of the full $240$-root shell into ten selected
$D_4$ subsystems:
\[
  \boxed{
  \text{obstruction regulus}
  \longrightarrow
  96\text{ exposed }E_8\text{ roots}
  \longrightarrow
  96+144=240\text{ roots uniquely}.
  }
\]
Each of the forty-five packets lies on exactly three such completion charts,
and each chart contains ten of the $270$ obstruction edges.  For every
obstruction edge, all four cross choices of one $D_4$ half from each packet
span rank eight; the verifier performs $270\cdot4=1080$ exact rank checks.

This theorem is deliberately finite-combinatorial.  It does not assert that an
$E_8$ completion removes a physical blocking obstruction, nor does it assign
dynamics to the root system.  It says that the independently certified
Holotrade obstruction edge is a canonical address of one selected $E_8$
ten-$D_4$ completion chart.

\subsection{The Schlaefli scheme on the twenty-seven completion charts}
\label{sec:20260901-e8-partition-schlafli}

The twenty-seven completion charts themselves recover the cubic-surface
incidence graph.  Regard each chart as a partition of the same $240$ roots into
ten selected $D_4$ blocks, and join two charts when they have a common
$D_4\oplus D_4$ packet block.  Any two charts have either zero or two common
$D_4$ blocks.  There are $216$ pairs of the first kind and $135$ of the second,
and the overlap graph is exactly
\[
  \boxed{\operatorname{SRG}(27,10,1,5)}.
\]
This is the line-intersection graph of the twenty-seven lines on a smooth cubic
surface; its complement is the usual Schlaefli graph
$\operatorname{SRG}(27,16,10,8)$.

The packet layer is recovered without extra labels.  Every one of the
forty-five selected $D_4\oplus D_4$ packets lies in exactly three completion
charts, those three charts form a triangle, and the overlap graph has exactly
forty-five triangles.  Conversely every triangle comes from exactly one
packet.  Thus, inside the selected $E_8$ partition catalogue,
\[
  \boxed{27\text{ cubic lines}
  \longleftrightarrow27\text{ ten-}D_4\text{ partitions}}
\]
and
\[
  \boxed{45\text{ tritangent planes}
  \longleftrightarrow45\text{ shared }D_4\oplus D_4\text{ packet blocks}}.
\]
Incidence is literal containment of partition blocks.

All twenty-seven charts cover the same $240$ roots, so the phrase ``two charts
intersect in $48$ roots'' would be misleading.  Their underlying root sets are
identical; the invariant proved here is the number of common \emph{partition
blocks}: two for adjacent charts and zero for nonadjacent charts.
